The rational relative-to-absolute degree-four map is an isomorphism
because $H^3(T^+_\Phi;\mathbb Q)=H^4(T^+_\Phi;\mathbb Q)=0$.
For each chosen spin structure, we use the classical Eells--Kuiper
invariant in \cite[Section~1.4, equations~(2)--(3), p.~5 of
arXiv:1406.2226v3]{mapsCrowleyNordstrom19}. Its source normalization has
$p_W=p_1(W)/2$ and values in $\mathbb Q/28\mathbb Z$; division by $28$
gives exactly
\[
 \mu=\frac{p_1^2-4\sigma}{896}\pmod{\mathbb Z}
\]
in our convention. This cited invariant applies to closed spin
seven-manifolds with torsion $p_W|_{\partial W}$, and that class is zero
here. Orientation reversal changes its sign.
